\documentclass[10pt,a4paper]{amsart}
\usepackage{amsmath,amssymb,amsthm,graphicx}
\usepackage{verbatim}
\newenvironment{reference}{\comment}{\endcomment}
%\newenvironment{reference}{}{}
\newenvironment{slogan}{\comment}{\endcomment}
\newenvironment{history}{\comment}{\endcomment}
% For commutative diagrams we use Xy-pic
\usepackage[all]{xy}

% We use 2cell for 2-commutative diagrams.
\xyoption{2cell}
% Theorem environments.
%
\theoremstyle{plain}
\newtheorem{theorem}[subsection]{Theorem}
\newtheorem{proposition}[subsection]{Proposition}
\newtheorem{lemma}[subsection]{Lemma}

\theoremstyle{definition}
\newtheorem{definition}[subsection]{Definition}
\newtheorem{example}[subsection]{Example}
\newtheorem{exercise}[subsection]{Exercise}
\newtheorem{situation}[subsection]{Situation}

\theoremstyle{remark}
\newtheorem{remark}[subsection]{Remark}
\newtheorem{remarks}[subsection]{Remarks}

\numberwithin{equation}{subsection}

% Macros
%
\def\lim{\mathop{\mathrm{lim}}\nolimits}
\def\colim{\mathop{\mathrm{colim}}\nolimits}
\def\Spec{\mathop{\mathrm{Spec}}}
\def\Hom{\mathop{\mathrm{Hom}}\nolimits}
\def\Ext{\mathop{\mathrm{Ext}}\nolimits}
\def\SheafHom{\mathop{\mathcal{H}\!\mathit{om}}\nolimits}
\def\SheafExt{\mathop{\mathcal{E}\!\mathit{xt}}\nolimits}
\def\Sch{\mathit{Sch}}
\def\Mor{\mathop{\mathrm{Mor}}\nolimits}
\def\Ob{\mathop{\mathrm{Ob}}\nolimits}
\def\Sh{\mathop{\mathit{Sh}}\nolimits}
\def\NL{\mathop{N\!L}\nolimits}
\def\CH{\mathop{\mathrm{CH}}\nolimits}
\def\proetale{{pro\text{-}\acute{e}tale}}
\def\etale{{\acute{e}tale}}
\def\QCoh{\mathit{QCoh}}
\def\Ker{\mathop{\mathrm{Ker}}}
\def\Im{\mathop{\mathrm{Im}}}
\def\Coker{\mathop{\mathrm{Coker}}}
\def\Coim{\mathop{\mathrm{Coim}}}

% Boxtimes
%
\DeclareMathSymbol{\boxtimes}{\mathbin}{AMSa}{"02}

%
% Macros for moduli stacks/spaces
%
\def\QCohstack{\mathcal{QC}\!\mathit{oh}}
\def\Cohstack{\mathcal{C}\!\mathit{oh}}
\def\Spacesstack{\mathcal{S}\!\mathit{paces}}
\def\Quotfunctor{\mathrm{Quot}}
\def\Hilbfunctor{\mathrm{Hilb}}
\def\Curvesstack{\mathcal{C}\!\mathit{urves}}
\def\Polarizedstack{\mathcal{P}\!\mathit{olarized}}
\def\Complexesstack{\mathcal{C}\!\mathit{omplexes}}
% \Pic is the operator that assigns to X its picard group, usage \Pic(X)
% \Picardstack_{X/B} denotes the Picard stack of X over B
% \Picardfunctor_{X/B} denotes the Picard functor of X over B
\def\Pic{\mathop{\mathrm{Pic}}\nolimits}
\def\Picardstack{\mathcal{P}\!\mathit{ic}}
\def\Picardfunctor{\mathrm{Pic}}
\def\Deformationcategory{\mathcal{D}\!\mathit{ef}}
\usepackage{cleveref}
\crefname{originallemma}{original lemma}{original lemmas}
\Crefname{originallemma}{Original lemma}{Original lemmas}
\makeatletter
\newlabel{lemma-essential-extension}{{lemma-essential-extension}{source}{}{sourcepointer-1}{}}
\newlabel{lemma-essential-extension@cref}{{[originallemma][0][]lemma-essential-extension}{[1][0][]source}{}{}{}}
\makeatother
\crefname{originallemma}{original lemma}{original lemmas}
\Crefname{originallemma}{Original lemma}{Original lemmas}
\makeatletter
\newlabel{lemma-injective-epimorphism}{{lemma-injective-epimorphism}{source}{}{sourcepointer-2}{}}
\newlabel{lemma-injective-epimorphism@cref}{{[originallemma][0][]lemma-injective-epimorphism}{[1][0][]source}{}{}{}}
\makeatother
\crefname{originallemma}{original lemma}{original lemmas}
\Crefname{originallemma}{Original lemma}{Original lemmas}
\makeatletter
\newlabel{lemma-injective-hull-unique}{{lemma-injective-hull-unique}{source}{}{sourcepointer-3}{}}
\newlabel{lemma-injective-hull-unique@cref}{{[originallemma][0][]lemma-injective-hull-unique}{[1][0][]source}{}{}{}}
\makeatother
\title{The injective hull of a domain modulo a prime is indecomposable and localizes}
\author{The Stacks Project Authors}
\date{stacks-project:dualizing.tex:local-snapshot-20261012; GNU FDL 1.2 or later}
\begin{document}
\maketitle
\noindent\textit{Editorial note.} All delivered mathematical text is the author's own, copied from the pinned Stacks Project chapter source sources/stacks-dualizing/pinned-original-source.tex; no statement or proof step was rewritten, reordered or removed.\par\medskip
\noindent\textit{Editorial note.} Preamble provenance: the Stacks Project keeps one single shared preamble for every chapter; this card takes it directly from the staged sources/stacks-dualizing/preamble.tex as three verbatim spans (lines 8-12, 14-18 and 154-217), so no derived file was written and nothing was added or redefined. The wrapper class and the base packages are supplied by the assembler.\par\medskip
\noindent\textit{Editorial note.} Bibliography provenance: neither the counted statement nor its proof contains a citation (this lane rescans the whole span before writing the card), so the card carries no bibliography and no entry was reconstructed.\par\medskip
\noindent\textit{Editorial note.} Selection provenance: this card was opened under the file-based control contract (CONTROL.md), with the mechanical pre-card gate run on its statement and proof interval before the card was written.\par\medskip
\noindent\textit{Allowed standard prerequisites (editorial).}\begin{itemize}
\item The essentialness of the extension of the residue class domain into its fraction field. Invoked inside the counted proof (source lines 826-828 and 839-840). Identifies the injective hull of the quotient with that of the residue field and detects nonzero intersection with submodules.
\item The uniqueness of injective hulls and the criterion for an inclusion of injective hulls to be an isomorphism. Invoked inside the counted proof (source lines 828-829, 831-832, 837-838 and 840-841). Identifies hulls computed in different ways and upgrades inclusions of injective submodules to isomorphisms.
\item The injectivity of a module over a localization acted on invertibly. Invoked inside the counted proof (source line 834). Shows the hull is injective as a module over the local ring at the prime.
\end{itemize}\medskip
\ifdefined\hypertarget\hypertarget{sourcepointer-1}{}\fi
\noindent\textit{Source reference.} Original lemma lemma-essential-extension: sources/stacks-dualizing/pinned-original-source.tex, source lines 371--371. Navigation pointer only: the counted text refers to the Stacks item lemma-essential-extension, which is an existing labeled result of this same chapter and is not part of this card. The pointer carries the target's own label as its number and its definition line as its span; its content is not reproduced here.\par\medskip
\ifdefined\hypertarget\hypertarget{sourcepointer-2}{}\fi
\noindent\textit{Source reference.} Original lemma lemma-injective-epimorphism: sources/stacks-dualizing/pinned-original-source.tex, source lines 422--422. Navigation pointer only: the counted text refers to the Stacks item lemma-injective-epimorphism, which is an existing labeled result of this same chapter and is not part of this card. The pointer carries the target's own label as its number and its definition line as its span; its content is not reproduced here.\par\medskip
\ifdefined\hypertarget\hypertarget{sourcepointer-3}{}\fi
\noindent\textit{Source reference.} Original lemma lemma-injective-hull-unique: sources/stacks-dualizing/pinned-original-source.tex, source lines 719--719. Navigation pointer only: the counted text refers to the Stacks item lemma-injective-hull-unique, which is an existing labeled result of this same chapter and is not part of this card. The pointer carries the target's own label as its number and its definition line as its span; its content is not reproduced here.\par\medskip
\begin{lemma}
\label{lemma-injective-hull-indecomposable}
Let $\mathfrak p \subset R$ be a prime of a ring $R$.
Let $E$ be the injective hull of $R/\mathfrak p$. Then
\begin{enumerate}
\item $E$ is indecomposable,
\item $E$ is the injective hull of $\kappa(\mathfrak p)$,
\item $E$ is the injective hull of $\kappa(\mathfrak p)$
over the ring $R_\mathfrak p$.
\end{enumerate}
\end{lemma}

\begin{proof}
By Lemma \ref{lemma-essential-extension} the inclusion
$R/\mathfrak p \subset \kappa(\mathfrak p)$ is an essential
extension. Then Lemma \ref{lemma-injective-hull-unique}
shows (2) holds. For $f \in R$, $f \not \in \mathfrak p$
the map $f : \kappa(\mathfrak p) \to \kappa(\mathfrak p)$ is an isomorphism
hence the map $f : E \to E$ is an isomorphism,
see Lemma \ref{lemma-injective-hull-unique}.
Thus $E$ is an $R_\mathfrak p$-module. It is injective
as an $R_\mathfrak p$-module by Lemma \ref{lemma-injective-epimorphism}.
Finally, let $E' \subset E$ be a nonzero injective $R$-submodule.
Then $J = (R/\mathfrak p) \cap E'$ is nonzero. After shrinking $E'$
we may assume that $E'$ is the injective hull of $J$ (see
Lemma \ref{lemma-injective-hull-unique} for example).
Observe that $R/\mathfrak p$ is an essential extension of $J$ for example by
Lemma \ref{lemma-essential-extension}. Hence $E' \to E$
is an isomorphism by Lemma \ref{lemma-injective-hull-unique} part (3).
Hence $E$ is indecomposable.
\end{proof}

\end{document}
